%=====================================================================
\section{Empirical process}
\label{sec:empirical}

The section describes the experimental procedure in sufficient detail for the
reproduction of the reported results. The section first argues the choice
of ensemble method, and then states the partitioning of the data, the resampling protocol, the tuning design of every model and
the justification of every candidate value, the protocol of the final
evaluation, and the statistical procedure used to compare the models.

\subsection{Choice of ensemble method}
\label{sec:choice}

The comparison was designed to reflect the effect of ensembling rather than
the effect of a different hypothesis space. The \gls{rf} uses the
classification tree of the preceding study as the base learner, so a
difference between the two models measures what aggregation and
randomisation add to the same learner. A dissimilar ensemble would confound
the two effects. The \gls{rf} is also a strong default for tabular data:
Fern\'{a}ndez-Delgado \textit{et al.}~\cite{fernandez2014} evaluated 179
classifiers on 121 datasets and found the \gls{rf} family to be the most
accurate overall.

Three properties of the data informed the choice among tree ensembles. The
first property is the label conflict described in
Section~\ref{sec:dataset}. From the perspective of any classifier, records
with identical features and different categories are label noise. Boosting
increases the weight of every misclassified instance, and would therefore
concentrate on records that no model separates.
Dietterich~\cite{dietterich2000b} showed that AdaBoost degraded markedly
under class label noise, whereas bagging and randomised trees did not. The
second property is the instability of the single tree under perturbation of
the training data, which is the condition under which averaging over
bootstrap samples succeeds~\cite{breiman1996}. The third property is the
strong correlation between features reported in Section~\ref{sec:dataset}.
Correlated features cause bagged trees to select similar splits, and the
random feature subset of the \gls{rf} offers a direct control over the
resulting correlation between trees, as expressed by
Equation~(\ref{eq:bound}). Stacking was not considered, because stacking
requires heterogeneous members, and the \gls{knn} model was both the weakest
and the most expensive candidate member.

\subsection{Partitioning and resampling}
\label{sec:partition}

The deduplicated dataset was partitioned once, by stratified random
sampling, into a development partition of \num{130196} records
($80\,\%$) and a generalisation set of \num{32549} records ($20\,\%$).
Every tuning decision was based on the development partition alone.

Within the development partition, every configuration was evaluated by
stratified five-fold cross-validation. Stratification preserves the class
proportions within every fold, which is necessary for a dataset in which an
unstratified partition can omit a rare class~\cite{kohavi1995}. Five folds
were preferred to ten folds for two reasons. The worms category comprises
137 development records, which yields approximately 27 records per test
fold at five folds but only 14 at ten folds, so that the per-fold recall of
the category would rest on few instances. Every model was evaluated on the same five folds, so that every
comparison during tuning was paired.

The primary performance measure was $\bar{F}$ of
Equation~(\ref{eq:macro}). The normal and generic categories jointly account
for $57.3\,\%$ of the records, so a classifier that serves only the largest
categories attains a high value of $\eta$ while failing at the purpose of
an \gls{ids}. The measures $\bar{r}$ and $\eta$ were recorded in support. The
configuration of greatest mean $\bar{F}$ over the five folds was selected
at every tuning stage.

\subsection{Tuning of the nearest neighbour model}
\label{sec:tune-knn}

The seven pre-processing variants of Section~\ref{sec:prep-knn} were each
combined with nine values of $k$, drawn from
$\{1, 3, 5, 7, 9, 11, 15, 21, 31\}$, and both weighting schemes, for a total
of 126 configurations. Odd values of $k$ avoid tied votes between two
classes. The values are dense around five, the optimum of the preceding
study, and are spaced approximately geometrically above eleven, because the
decision boundary changes progressively less per additional neighbour as
$k$ grows.

\subsection{Tuning of the classification tree}
\label{sec:tune-tree}

The tree was tuned in three stages. The first stage evaluated 72
configurations, formed by the combination of both impurity measures, six
depth limits drawn from $\{8, 12, 16, 20, 24, \infty\}$, three minimum leaf
sizes drawn from $\{1, 5, 20\}$, and two weighting schemes, namely uniform
weights and weights in inverse proportion to class frequency. The depth
limits bracket the optimum of 16 found in the preceding study at a regular
step of four, and the unlimited depth serves as the unpruned reference. The
leaf sizes are spaced approximately geometrically, because the effect of a
leaf size limit scales with the logarithm of the limit. The second stage
applied eight values of $\alpha$ in Equation~(\ref{eq:ccp}), drawn from
$\{0, 1, 2, 5, 10, 20, 50, 100\} \times 10^{-5}$, to the best configuration
of the first stage. The values follow a one, two, five progression that
brackets the optimum of $10^{-4}$ of the preceding study. The third stage
evaluated the four encodings of Section~\ref{sec:prep-tree} under the best
configuration of the second stage.

\subsection{Tuning of the random forest}
\label{sec:tune-rf}

The forest was tuned in three stages, each of which addresses specific
research questions of Section~\ref{sec:intro}. A joint search over every
control parameter was computationally infeasible on the available hardware,
so the stages form a coordinate-wise search.

The first stage addressed the third and the fifth research questions. Eight
values of $q$, drawn from $\{1, 2, 4, 6, 10, 16, 24, 42\}$, were combined with
three weighting schemes: uniform weights, weights computed once from the
training partition, and weights recomputed for every bootstrap sample. The
values of $q$ span the full range from one feature per split, which
randomises the choice of split feature completely, to all 42 features, which
reduces the \gls{rf} to bagging. The spacing is approximately geometric, and
the value of six corresponds to $\sqrt{p}$, the default recommended by
Breiman~\cite{breiman2001} for classification. Trees were grown without depth
or leaf size limits, following the convention of Section~\ref{sec:bg-rf}.
Every configuration used 100 trees, and the strength and the pairwise
agreement of the members were recorded.

The second stage addressed the fourth research question. Under the best
configuration of the first stage, six minimum leaf sizes drawn from
$\{1, 2, 5, 10, 20, 50\}$ and six depth limits drawn from
$\{8, 12, 16, 20, 24, 32\}$ were evaluated one at a time. The depth limits
include the optimum of the single tree, so that the effect of the limit
that the single tree required was observed directly. The mean $\bar{F}$ on
the training partitions was recorded alongside the test partitions, so that
the gap between the two quantified overfitting.

The third stage addressed the second research question. For the best
configuration of the second stage, and for the bagged equivalent with $q$
equal to $p$, $\bar{F}$ was recorded for every value of $\tau$ from one to
300. The number of trees for the final evaluation was fixed beforehand as
the smallest value in $\{25, 50, 100, 150, 200, 250, 300\}$ whose mean
$\bar{F}$ lies within 0.001 of the value at 300 trees. A final check
compared the integer and indicator encodings under the selected
configuration, with $q$ held at the same fraction of the available features.

\subsection{Final evaluation}
\label{sec:final-protocol}

Four models were evaluated on the generalisation set: the \gls{knn} model,
the classification tree, the bagged ensemble and the \gls{rf}. Every model
was fitted to the full development partition under the selected
configuration. The bagged ensemble shares every control parameter of the
\gls{rf} except $q$, so the comparison of the two ensembles isolates the
effect of random feature selection, and the comparison of the bagged
ensemble with the tree isolates the effect of bootstrap aggregation.

The tree-based models depend on a random seed, through bootstrap sampling,
through the random feature subsets, and through the resolution of ties
between splits of equal impurity reduction. Each tree-based model was
therefore fitted under ten seeds, and the mean and the standard deviation
over the seeds are reported. The \gls{knn} model is deterministic and was
evaluated once. Training time, prediction time, per-class measures and the
confusion matrix were recorded for every model.

\subsection{Statistical comparison}
\label{sec:stats-protocol}

Five pairwise comparisons were planned: the \gls{rf} against each of the
other three models, and the bagged ensemble against the tree. Each
comparison was performed on the predictions of the models fitted under the
first seed. The difference in $\bar{F}$ was assessed with \num{2000}
paired bootstrap resamples of the generalisation set, and the difference in
accuracy with the McNemar test of Equation~(\ref{eq:mcnemar}). The
$p$-values of each test family were adjusted by the procedure of
Holm~\cite{holm1979}, and a significance level of 0.05 was applied.

\subsection{Irreducible error}
\label{sec:ceiling-protocol}

The records that share identical feature values with a record of a different
category cannot all be classified correctly by any deterministic classifier.
An empirical ceiling was computed by grouping the deduplicated records by
their feature values and assigning to every group the most frequent
category of the group. The ceiling is optimistic, because the assignment
uses the labels that the assignment then scores, but the ceiling identifies
the categories for which the features are insufficient.
